
\subsection{\texorpdfstring{$A_\infty$}{A-infinity} 双模同态的形变}
\label{subsec:5-2-2}

第 3.6.1 与 3.7.3 小节分别定义了
滤过 $A_\infty$ 代数与滤过 $A_\infty$ 双模的形变。
本小节定义滤过 $A_\infty$ 代数同态
及滤过 $A_\infty$ 双模同态的形变。
为此，先引入滤过 $A_\infty$ 双模结构的\emph{拉回}。
设 $(D,\mathfrak n)$、$(D',\mathfrak n')$
是上一小节中以 $\Lambda_{0,\mathrm{nov}}$
为系数的滤过 $A_\infty$ 双模。

\begin{lemma}{5.2.6}\label{lem:5-2-6}
设 $(D,\mathfrak n)$ 是
$(C'_1,C'_0)$ 上的滤过 $A_\infty$ 双模，
$\mathfrak f^{(i)}:C_i\to C'_i$
（$i=0,1$）是滤过 $A_\infty$ 同态。
则 $D$ 上存在一个 $(C_1,C_0)$-滤过
$A_\infty$ 双模结构，记为
\[
\bigl((\mathfrak f^{(1)},\mathfrak f^{(0)})^*D,
      (\mathfrak f^{(1)},\mathfrak f^{(0)})^*\mathfrak n\bigr).
\]
此外，在 $(\mathfrak f^{(1)},\mathfrak f^{(0)})$
之上存在滤过 $A_\infty$ 双模同态
\[
\varphi:
(\mathfrak f^{(1)},\mathfrak f^{(0)})^*D
\longrightarrow D,
\qquad \varphi_{0,0}=\mathrm{id}.
\]
非滤过或弱滤过的版本，
以及含幺或同伦含幺的版本，也都成立。
\translatornote{原书实体页 286
在本引理的同态源端把
$(\mathfrak f^{(1)},\mathfrak f^{(0)})^*D$
误写成 $(\mathfrak f^{(1)},\mathfrak f^{(2)})^*D$；
文中只定义 $i=0,1$ 的同态。
见疑误 A-0269。}
\end{lemma}

\begin{remark}{5.2.7}\label{rem:5-2-7}
本处及下文所说的
“非滤过或弱滤过的版本，
以及／或者含幺或同伦含幺的版本”，
涵盖以下九种情形：
非滤过、弱滤过、滤过各自与
含幺、同伦含幺、非含幺组合所得的版本。
具体而言，它们是非滤过含幺、弱滤过含幺、
滤过含幺、非滤过同伦含幺、
弱滤过同伦含幺、滤过同伦含幺、
非滤过非含幺、弱滤过非含幺和滤过非含幺版本。
\end{remark}

\begin{definition}{5.2.8}\label{def:5-2-8}
称 $(\mathfrak f^{(1)},\mathfrak f^{(0)})^*D$
为双模 $D$ 沿
$(\mathfrak f^{(1)},\mathfrak f^{(0)})$
的\emph{拉回}。
\liftIndex{I321}
\end{definition}

\begin{proof}[引理 5.2.6 的证明]
定义
\[
\begin{aligned}
(\mathfrak f^{(1)},\mathfrak f^{(0)})^*\mathfrak n:
&\widehat B(C_1[1])\widehat\otimes D[1]
 \widehat\otimes\widehat B(C_0[1])\\
&\longrightarrow D[1]
\end{aligned}
\]
为
\[
\begin{aligned}
&(\mathfrak f^{(1)},\mathfrak f^{(0)})^*
\mathfrak n(\mathbf x,y,\mathbf z)\\
&\qquad=
\mathfrak n\bigl(
\widehat{\mathfrak f}^{(1)}(\mathbf x),
y,\widehat{\mathfrak f}^{(0)}(\mathbf z)
\bigr).
\end{aligned}
\]
右端的 $\mathfrak n$ 是 $D$ 上既定的
$(C'_1,C'_0)$-$A_\infty$ 双模结构。
容易验证，上式定义了滤过的
$(C_1,C_0)$-$A_\infty$ 双模结构。
再置 $\varphi_{k_1,k_0}=0$
（若 $(k_1,k_0)\ne(0,0)$），
而 $\varphi_{0,0}=\mathrm{id}$，
即可验证 $\varphi$ 是位于
$(\mathfrak f^{(1)},\mathfrak f^{(0)})$
之上的滤过 $A_\infty$ 双模同态。
\end{proof}

拉回运算具有函子性。
在引理 5.2.6 的情形下，
再设 $\mathfrak g^{(i)}:C'_i\to C''_i$
（$i=0,1$）是滤过 $A_\infty$ 同态，
$D'$ 是 $(C''_1,C''_0)$ 上的滤过 $A_\infty$ 双模。

\begin{lemma}{5.2.9}\label{lem:5-2-9}
位于一对滤过 $A_\infty$ 同态
$(\mathfrak g^{(1)},\mathfrak g^{(0)})$
之上的滤过 $A_\infty$ 双模同态
$\varphi:D\to D'$，诱导位于恒等同态之上的
滤过 $A_\infty$ 双模同态
\[
\begin{aligned}
&(\mathfrak f^{(1)},\mathfrak f^{(0)})^*\varphi:\
(\mathfrak f^{(1)},\mathfrak f^{(0)})^*D\\
&\quad\longrightarrow
(\mathfrak g^{(1)}\circ\mathfrak f^{(1)},
 \mathfrak g^{(0)}\circ\mathfrak f^{(0)})^*D'.
\end{aligned}
\]
非滤过或弱滤过的版本，
以及／或者含幺或同伦含幺的版本，
也都成立。
\end{lemma}

\begin{proof}
定义
\[
\begin{aligned}
&\bigl((\mathfrak f^{(1)},\mathfrak f^{(0)})^*
\varphi\bigr)(\mathbf x,y,\mathbf z)\\
&\qquad=
\varphi\bigl(
\widehat{\mathfrak f}^{(1)}(\mathbf x),
y,\widehat{\mathfrak f}^{(0)}(\mathbf z)
\bigr)
\end{aligned}
\]
即可。
\end{proof}

\begin{definition}{5.2.10}\label{def:5-2-10}
称 $(\mathfrak f^{(1)},\mathfrak f^{(0)})^*\varphi$
为同态 $\varphi$ 沿
$(\mathfrak f^{(1)},\mathfrak f^{(0)})$
的\emph{拉回}。
\liftIndex{I325}
\end{definition}
\subsubsection{余代数同构的定义}
\label{subsubsec:5-2-2-1}

现在定义滤过 $A_\infty$ 代数同态的形变。
如第 3.6.2 小节所见，
对满足 $b\equiv0\pmod{\Lambda_{0,\mathrm{nov}}^+}$
的任意 $b\in C[1]^0$，
滤过 $A_\infty$ 代数 $(C,\mathfrak m)$
可以形变，使得
\[
\mathfrak m_k^b(x_1,\ldots,x_k)
=\mathfrak m(e^b x_1e^b x_2\cdots e^b x_ke^b).
\]
参见定义 3.6.9 与命题 3.6.10。
与此对应，下面形变滤过 $A_\infty$ 代数同态；
为此先在 bar 复形 $\widehat B(C[1])$
上引入一个\emph{余代数同构}\liftIndex{I184}。

对满足 $b\equiv0\pmod{\Lambda_{0,\mathrm{nov}}^+}$
的 $b\in C[1]^0$，
定义 $\mathfrak t_{b,k}:B_k(C[1])\to C[1]$：
$\mathfrak t_{b,0}(1)=b$，
而 $k\ne0$ 时 $\mathfrak t_{b,k}=0$。
它诱导余代数同态
\[
\widehat{\mathfrak t}_b:
\widehat B(C[1])\longrightarrow\widehat B(C[1]).
\]
其在 $B_k(C[1])$ 上的限制
$\widehat{\mathfrak t}_{b,k}$ 为
\begin{equation}
\begin{aligned}
\widehat{\mathfrak t}_{b,k}(x_1,\ldots,x_k)
=e^b\otimes x_1\otimes e^b\otimes\cdots
\otimes e^b\otimes x_k\otimes e^b.
\end{aligned}
\tag{5.2.11}\label{eq:5-2-11}
\end{equation}
由于 $b\equiv0\pmod{\Lambda_{0,\mathrm{nov}}^+}$，
有 $\widehat{\mathfrak t}_b\equiv\mathrm{id}
\pmod{\Lambda_{0,\mathrm{nov}}^+}$，
而 $\widehat{\mathfrak t}_b$
不降低任一项的能量。
所以它保持能量滤过，
并在相对于能量滤过完备化的
$\widehat B(C[1])$ 中收敛。

\begin{lemma}{5.2.12}\label{lem:5-2-12}
\textup{(1)} $\widehat{\mathfrak t}_b$
是 $\widehat B(C[1])$ 上的余代数同构。

\textup{(2)} 此外，
$\widehat d\circ\widehat{\mathfrak t}_b
=\widehat{\mathfrak t}_b\circ\widehat d^b$。
因此它定义滤过 $A_\infty$ 同态，记为
\[
\mathfrak t_b:(C,\mathfrak m^b)
\longrightarrow(C,\mathfrak m).
\]
\end{lemma}

\begin{proof}
\textup{(1)} 单射性显然。
为证明满射性，
对任意
$\mathbf y=y_1\otimes\cdots\otimes y_\ell
\in B_\ell(C[1])$，
逐次求解方程
\begin{equation}
\widehat{\mathfrak t}_b(\mathbf x)=\mathbf y.
\tag{5.2.13}\label{eq:5-2-13}
\end{equation}
对非负整数 $K$，
若 $\mathbf x\in\widehat B(C[1])$ 满足
\[
\left.\widehat{\mathfrak t}_b(\mathbf x)
\right|_{B_{1\cdots(\ell+K)}(C[1])}
=\mathbf y,
\]
便称它为式~\eqref{eq:5-2-13}
的一个\emph{$K$-近似解}。
这里如式 (4.4.5) 记
\[
B_{1\cdots(\ell+K)}(C[1])
=\bigoplus_{k=1}^{\ell+K}B_k(C[1]);
\]
另记
\[
B_{\ell\cdots(\ell+K)}(C[1])
=\bigoplus_{k=\ell}^{\ell+K}B_k(C[1]).
\]

\begin{statement}{子引理}{5.2.14}\label{sublem:5-2-14}
对任意
$\mathbf y=y_1\otimes\cdots\otimes y_\ell
\in B_\ell(C[1])$，
存在一族整数
$\{a_{i_0,\ldots,i_\ell}\}_{i_j\geq0,\,\ell>0}$，
使得对每个 $K=0,1,2,\ldots$，
\[
\begin{aligned}
\mathbf x_{(\ell+K)}
=\sum_{i_0+\cdots+i_\ell\leq K}
&a_{i_0,\ldots,i_\ell}
b^{\otimes i_0}\otimes y_1\otimes\cdots\\
&\otimes y_\ell\otimes b^{\otimes i_\ell}
\in B_{\ell\cdots(\ell+K)}(C[1])
\end{aligned}
\]
是式~\eqref{eq:5-2-13}
的一个 $K$-近似解。
\end{statement}

\begin{proof}[子引理 5.2.14 的证明]
对 $K$ 作归纳。
先注意到
\[
\widehat{\mathfrak t}_{b,k}:
B_k(C[1])\longrightarrow
\widehat{\bigoplus}_{\ell\geq k}B_\ell(C[1]),
\]
且 $\pi_k\circ\widehat{\mathfrak t}_{b,k}
=\mathrm{id}$；
其中 $\pi_k:\widehat{\bigoplus}_{\ell\geq k}
B_\ell(C[1])\to B_k(C[1])$ 是投影。
当 $K=0$ 时，取 $a_{0,\ldots,0}=1$，
于是
$\mathbf x_{(\ell)}=y_1\otimes\cdots\otimes y_\ell$
显然是 $0$-近似解。
假设
\[
\begin{aligned}
\mathbf x_{(\ell+K)}
=\sum_{i_0+\cdots+i_\ell\leq K}
&a_{i_0,\ldots,i_\ell}
b^{\otimes i_0}\otimes y_1\otimes\cdots\\
&\otimes y_\ell\otimes b^{\otimes i_\ell}
\end{aligned}
\]
是 $K$-近似解。
由归纳假设，
$\widehat{\mathfrak t}_b(\mathbf x_{(\ell+K)})$
在 $B_{1\cdots(\ell+K)}(C[1])$ 上的分量
等于 $\mathbf y$；于是对某些整数
$a'_{i_0,\ldots,i_\ell}$，可写作
\[
\begin{aligned}
\widehat{\mathfrak t}_b(\mathbf x_{(\ell+K)})
=\mathbf y+
\sum_{i_0+\cdots+i_\ell>K}
&a'_{i_0,\ldots,i_\ell}
b^{\otimes i_0}\otimes y_1\otimes\cdots\\
&\otimes y_\ell\otimes b^{\otimes i_\ell}.
\end{aligned}
\]
令 $\operatorname{pr}_{\leq\ell+K+1}$
为向不超过 $\ell+K+1$ 的张量长度的投影。
定义\translatornote{原书实体页 288
的修正项直接使用未截断的
$\widehat{\mathfrak t}_b$，
它会产生任意高张量长度的项，
与子引理要求
$\mathbf x_{(\ell+K+1)}\in
B_{\ell\cdots(\ell+K+1)}(C[1])$
不符。此处只取下一阶所需的长度截断，
并统一相容式中的括号下标；
见疑误 A-0270。}
\[
\begin{aligned}
\mathbf x_{(\ell+K+1)}
=\mathbf x_{(\ell+K)}-
\sum_{i_0+\cdots+i_\ell=K+1}
&a'_{i_0,\ldots,i_\ell}
\operatorname{pr}_{\leq\ell+K+1}
\widehat{\mathfrak t}_b\bigl(
b^{\otimes i_0}\otimes y_1\otimes\cdots\\
&\otimes y_\ell\otimes b^{\otimes i_\ell}
\bigr).
\end{aligned}
\]
容易看出，$\mathbf x_{(\ell+K+1)}$
是 $(K+1)$-近似解，且满足
\[
\mathbf x_{(\ell+K+1)}
\equiv\mathbf x_{(\ell+K)}
\pmod{\widehat{\bigoplus}_{k\geq\ell+K+1}B_k(C[1])}.
\]
这就证明了子引理。
\end{proof}

上证最后一个关系表明，
直极限
\[
\mathbf x:=\lim_{K\to\infty}
\mathbf x_{(\ell+K)}
\]
存在。
又因 $b\equiv0\pmod{\Lambda_{0,\mathrm{nov}}^+}$，
此极限也在相对于能量滤过的完备化
$\widehat B(C[1])$ 中收敛，
从而解出式~\eqref{eq:5-2-13}。
引理 5.2.12 的第 \textup{(1)} 项得证。

\textup{(2)} 直接计算可得
$\widehat d\circ\widehat{\mathfrak t}_b
=\widehat{\mathfrak t}_b\circ\widehat d^{\,b}$。
而且如引理 5.2.12 前所述，
$\widehat{\mathfrak t}_b$ 保持能量滤过，
所以满足 (3.2.25.1)。
\end{proof}

\subsubsection{\texorpdfstring{$A_\infty$}{A-infinity} 代数同态的形变}
\label{subsubsec:5-2-2-2}

现在考虑滤过 $A_\infty$ 代数
$(C,\mathfrak m)$ 与 $(C',\mathfrak m')$ 之间
的滤过 $A_\infty$ 同态
$\mathfrak f:C\to C'$。
记其诱导的余代数同态为
$\widehat{\mathfrak f}:\widehat B(C[1])
\to\widehat B(C'[1])$（见式 (3.2.28)）。
由引理 3.6.37 的第 \textup{(1)} 项，
回顾对任意
$b\in C[1]^0$ 且
$b\equiv0\pmod{\Lambda_{0,\mathrm{nov}}^+}$，
已定义
\[
\begin{aligned}
\mathfrak f_*(b)
&=\mathfrak f(e^b)\\
&=\mathfrak f_0(1)+\mathfrak f_1(b)
+\mathfrak f_2(b,b)+\cdots
\in C'[1]^0,
\end{aligned}
\]
其中 $\mathfrak f=\sum_{k=0}^{\infty}\mathfrak f_k$。
注意 $\mathfrak f_*(b)\equiv0
\pmod{\Lambda_{0,\mathrm{nov}}^+}$，
所以还可在 $\widehat B(C'[1])$
上考虑余代数同构
$\widehat{\mathfrak t}_{\mathfrak f_*(b)}$。

\begin{statement}{定义--引理}{5.2.15}\label{deflem:5-2-15}
对满足上述条件的 $b\in C[1]^0$，定义
\[
\begin{aligned}
\widehat{\mathfrak f}^{b}
&=\widehat{\mathfrak t}_{\mathfrak f_*(b)}^{-1}
\circ\widehat{\mathfrak f}
\circ\widehat{\mathfrak t}_b:\\
&\widehat B(C[1])\longrightarrow\widehat B(C'[1]).
\end{aligned}
\]
它使图式 5.2.1 交换。
{\renewcommand{\thefigure}{5.2.1}
\begin{center}
\[
\begin{array}{ccc}
(\widehat B(C[1]),\widehat d)
&\xrightarrow{\widehat{\mathfrak f}}
&(\widehat B(C'[1]),\widehat d')\\[0.5ex]
\scriptstyle{\widehat{\mathfrak t}_b}\Big\uparrow
&&\Big\downarrow\scriptstyle{
\widehat{\mathfrak t}_{\mathfrak f_*(b)}^{-1}}\\[0.5ex]
(\widehat B(C[1]),\widehat d^b)
&\xrightarrow{\widehat{\mathfrak f}^{\,b}}
&\bigl(\widehat B(C'[1]),
\widehat d'^{\,\mathfrak f_*(b)}\bigr)
\end{array}
\]
\refstepcounter{figure}
{\normalfont\bfseries 图式~\thefigure}\label{diag:5-2-1}
\end{center}}
$\widehat{\mathfrak f}^{b}$ 显然是余代数同态。
由引理 5.2.12 的第 \textup{(2)} 项，
有
\[
\widehat{\mathfrak f}^{b}\circ\widehat d^b
=\widehat d'^{\,\mathfrak f_*(b)}
\circ\widehat{\mathfrak f}^{b}.
\]
此外，
$\widehat{\mathfrak t}_{\mathfrak f_*(b)}^{-1}$
也模 $\Lambda_{0,\mathrm{nov}}^+$ 同余于恒等映射；
由子引理 5.2.14 的构造，
其正能量部分的能量上升。
故 $\widehat{\mathfrak f}^{b}$
满足 (3.2.27.1)，定义滤过 $A_\infty$ 同态
\[
\mathfrak f^b:(C,\mathfrak m^b)
\longrightarrow
(C',\mathfrak m'^{\,\mathfrak f_*(b)}).
\]
称 $\mathfrak f^b$ 为
$\mathfrak f:C\to C'$ 沿 $b$ 的\emph{形变}\liftIndex{I326}\liftIndex{I330}。
\end{statement}

下面的引理在论证中具有基础作用。

\begin{lemma}{5.2.16}\label{lem:5-2-16}
形变后的滤过 $A_\infty$ 代数同态
$\mathfrak f^b$ 总是严格的，
即 $\mathfrak f^b_0=0$。
\end{lemma}

\begin{proof}
只须证明 $\mathfrak f^b_0(1)=0$。
注意
\[
\widehat{\mathfrak f}(e^b)
=e^{\mathfrak f_*(b)},
\qquad
\widehat{\mathfrak f}^{b}(1)
=1+\mathfrak f^b_0(1)
+\mathfrak f^b_0(1)\otimes\mathfrak f^b_0(1)
+\cdots
=e^{\mathfrak f^b_0(1)}.
\]
依 $\widehat{\mathfrak t}_b$ 的定义，
\[
\begin{aligned}
e^{\mathfrak f^b_0(1)}
&=\widehat{\mathfrak f}^{b}(1)
=\widehat{\mathfrak t}_{\mathfrak f_*(b)}^{-1}
\circ\widehat{\mathfrak f}(e^b)\\
&=\widehat{\mathfrak t}_{\mathfrak f_*(b)}^{-1}
\bigl(e^{\mathfrak f_*(b)}\bigr)=1.
\end{aligned}
\]
故 $\mathfrak f^b_0(1)=0$。
\end{proof}

含幺和同伦含幺的情形也成立：

\begin{lemma}{5.2.17}\label{lem:5-2-17}
设 $(C,\mathfrak m)$、$(C',\mathfrak m')$
是含幺（相应地，同伦含幺）的
滤过 $A_\infty$ 代数，
$\mathfrak f:(C,\mathfrak m)\to(C',\mathfrak m')$
是含幺（相应地，同伦含幺）的
滤过 $A_\infty$ 同态。
则对任意 $b\in C[1]^0$ 且
$b\equiv0\pmod{\Lambda_{0,\mathrm{nov}}^+}$，
形变后的同态
\[
\mathfrak f^b:(C,\mathfrak m^b)
\longrightarrow
(C',\mathfrak m'^{\,\mathfrak f_*(b)})
\]
仍然含幺（相应地，同伦含幺）；
两种情形下它也都是严格的。
\end{lemma}

\subsubsection{\texorpdfstring{$A_\infty$}{A-infinity} 双模同态的形变}
\label{subsubsec:5-2-2-3}

接着定义滤过 $A_\infty$ 双模同态的形变；
其方式与刚讨论的滤过 $A_\infty$ 代数同态情形相似。
设 $\mathfrak f^{(i)}:C_i\to C'_i$
（$i=0,1$）是滤过 $A_\infty$ 代数同态，
$(D,\mathfrak n)$ 与 $(D',\mathfrak n')$
分别是 $(C_1,C_0)$ 与 $(C'_1,C'_0)$ 上的
滤过 $A_\infty$ 双模。
取位于
$(\mathfrak f^{(1)},\mathfrak f^{(0)})$
之上的滤过 $A_\infty$ 双模同态
$\varphi:(D,\mathfrak n)\to(D',\mathfrak n')$，
并按定义 3.7.7 记其诱导同态为
\[
\begin{aligned}
\widehat\varphi:
&\widehat B(C_1[1])\widehat\otimes D[1]
 \widehat\otimes\widehat B(C_0[1])\\
&\longrightarrow
\widehat B(C'_1[1])\widehat\otimes D'[1]
 \widehat\otimes\widehat B(C'_0[1]).
\end{aligned}
\]
记 $\widehat d,\widehat d'$
分别为式 (3.7.3) 中由
$\mathfrak n,\mathfrak n'$ 诱导的余导子。

对满足
$b_i\equiv0\pmod{\Lambda_{0,\mathrm{nov}}^+}$
的 $b_i\in C_i[1]^0$，
定义--引理 3.7.13 已给出形变双模
$(D,{}^{b_1}\widehat d^{b_0})$。
同样，$D'$ 沿
$\mathfrak f^{(i)}_*(b_i)\in C'_i[1]^0$
形变为
\[
\bigl(D',{}^{\mathfrak f^{(1)}_*(b_1)}
\widehat d'^{\,\mathfrak f^{(0)}_*(b_0)}\bigr),
\]
其左右分别承载
$(C'_1,\mathfrak m'^{\,\mathfrak f^{(1)}_*(b_1)})$
与
$(C'_0,\mathfrak m'^{\,\mathfrak f^{(0)}_*(b_0)})$
的作用。
另一方面，定义--引理 5.2.15 把
$\mathfrak f^{(i)}$ 沿 $b_i$ 形变为
\[
\bigl(\mathfrak f^{(i)}\bigr)^{b_i}:
(C_i,\mathfrak m^{b_i})
\longrightarrow
(C'_i,\mathfrak m'^{\,\mathfrak f^{(i)}_*(b_i)}),
\]
由引理 5.2.16，它们均为严格同态。

\begin{statement}{定义--引理}{5.2.18}\label{deflem:5-2-18}
对上述 $b_i\in C_i[1]^0$，定义
\[
\begin{aligned}
{}^{b_1}\widehat\varphi^{b_0}
={}&\bigl(
\widehat{\mathfrak t}_{\mathfrak f^{(1)}_*(b_1)}^{-1}
\otimes\mathrm{id}\otimes
\widehat{\mathfrak t}_{\mathfrak f^{(0)}_*(b_0)}^{-1}
\bigr)\\
&\quad\circ\widehat\varphi\circ
\bigl(\widehat{\mathfrak t}_{b_1}
\otimes\mathrm{id}\otimes
\widehat{\mathfrak t}_{b_0}\bigr).
\end{aligned}
\]
于是图式 5.2.2 交换。
为使图中的结点紧凑，记
\[
\begin{aligned}
\mathcal B_D
&=\widehat B(C_1[1])\widehat\otimes D[1]
\widehat\otimes\widehat B(C_0[1]),\\
\mathcal B'_{D'}
&=\widehat B(C'_1[1])\widehat\otimes D'[1]
\widehat\otimes\widehat B(C'_0[1]).
\end{aligned}
\]
{\renewcommand{\thefigure}{5.2.2}
\begin{center}
\resizebox{\linewidth}{!}{$\displaystyle
\begin{array}{ccc}
(\mathcal B_D,\widehat d)
&\xrightarrow{\widehat\varphi}
&(\mathcal B'_{D'},\widehat d')\\[1ex]
\scriptstyle{\widehat{\mathfrak t}_{b_1}
\otimes\mathrm{id}\otimes\widehat{\mathfrak t}_{b_0}}
\Big\uparrow
&&\Big\downarrow
\scriptstyle{\widehat{\mathfrak t}_{\mathfrak f^{(1)}_*(b_1)}^{-1}
\otimes\mathrm{id}\otimes
\widehat{\mathfrak t}_{\mathfrak f^{(0)}_*(b_0)}^{-1}}\\[1ex]
(\mathcal B_D,{}^{b_1}\widehat d^{b_0})
&\xrightarrow{{}^{b_1}\widehat\varphi^{b_0}}
&(\mathcal B'_{D'},
{}^{\mathfrak f^{(1)}_*(b_1)}
\widehat d'^{\,\mathfrak f^{(0)}_*(b_0)})
\end{array}$}
\refstepcounter{figure}
{\normalfont\bfseries 图式~\thefigure}\label{diag:5-2-2}
\end{center}}
因此 ${}^{b_1}\widehat\varphi^{b_0}$
在严格 $A_\infty$ 同态
$\bigl((\mathfrak f^{(1)})^{b_1},
(\mathfrak f^{(0)})^{b_0}\bigr)$
之上定义滤过 $A_\infty$ 双模同态
\[
{}^{b_1}\varphi^{b_0}:
(D,{}^{b_1}\widehat d^{b_0})
\longrightarrow
\bigl(D',{}^{\mathfrak f^{(1)}_*(b_1)}
\widehat d'^{\,\mathfrak f^{(0)}_*(b_0)}\bigr).
\]
含幺及同伦含幺的版本也成立；
弱滤过以及／或者含幺、同伦含幺的版本亦然。
称 ${}^{b_1}\varphi^{b_0}$ 为
$\varphi$ 沿 $b_1,b_0$ 的\emph{形变}\liftIndex{I322}\liftIndex{I331}。
\end{statement}

\begin{proof}
上述结论直接由定义检验，证明从略。
\end{proof}

\subsubsection{无障碍情形}
\label{subsubsec:5-2-2-4}

现假设 $C_i$（$i=0,1$）无障碍，
并取定界余链 $b_i\in\mathcal M(C_i)$。
若 $\mathfrak f^{(i)}:C_i\to C'_i$
为 $A_\infty$ 代数同态，
则引理 3.6.36 表明
$\mathfrak f^{(i)}_*(b_i)$ 也是定界余链。
故 $\mathfrak f^{(i)}_*(b_i)
\in\mathcal M(C'_i)\ne\varnothing$，
$C'_i$ 也无障碍。
由引理 3.7.14，存在复形
$(D,\delta_{b_1,b_0})$ 与
\[
\bigl(D',\delta_{\mathfrak f^{(1)}_*(b_1),
\mathfrak f^{(0)}_*(b_0)}\bigr).
\]
于是有下述结论。

\begin{lemma}{5.2.19}\label{lem:5-2-19}
在定义--引理 5.2.18 的情形下，
对任意定界余链
$b_i\in\mathcal M(C_i)$，
位于
$\bigl((\mathfrak f^{(1)})^{b_1},
(\mathfrak f^{(0)})^{b_0}\bigr)$
之上的形变 $A_\infty$ 双模同态
${}^{b_1}\varphi^{b_0}$
诱导上链映射
\[
{}^{b_1}\varphi^{b_0}:
(D,\delta_{b_1,b_0})
\longrightarrow
\bigl(D',\delta_{\mathfrak f^{(1)}_*(b_1),
\mathfrak f^{(0)}_*(b_0)}\bigr).
\]
这里对 $y\in D$，
\[
\delta_{b_1,b_0}(y)
={}^{b_1}\mathfrak n^{b_0}(y)
=\mathfrak n(e^{b_1}ye^{b_0});
\]
参见定义--引理 3.7.13 的第 \textup{(2)} 项。
若 $C_i$ 含幺且弱无障碍，
则对任意满足
$\mathfrak{PO}(b_1)=\mathfrak{PO}(b_0)$
的 $b_i\in\mathcal M_{\mathrm{weak}}(C_i)$，
上述诱导映射仍为上链映射。
弱滤过以及／或者同伦含幺的版本也成立。
\end{lemma}

\begin{proof}
由于
$\mathfrak m_0^{b_i}
=\mathfrak m_0'^{\,\mathfrak f^{(i)}_*(b_i)}=0$，
且由引理 5.2.16，
$(\mathfrak f^{(i)})^{b_i}$ 是严格同态，
只须证明下列一般事实。

\begin{statement}{子引理}{5.2.20}\label{sublem:5-2-20}
\textup{(1)} 设 $(D,\mathfrak n)$ 是
$\bigl((C_1,\mathfrak m^{(1)}),
(C_0,\mathfrak m^{(0)})\bigr)$
上的滤过 $A_\infty$ 双模。
若对 $i=0,1$ 均有
$\mathfrak m_0^{(i)}=0$，
则 $\mathfrak n_{0,0}\circ\mathfrak n_{0,0}=0$。

\textup{(2)} 进一步设 $(C'_i,\mathfrak m'^{(i)})$
为满足 $\mathfrak m_0'^{(i)}=0$ 的
滤过 $A_\infty$ 代数，
$(D',\mathfrak n')$ 是
$\bigl((C'_1,\mathfrak m'^{(1)}),
(C'_0,\mathfrak m'^{(0)})\bigr)$
上的滤过 $A_\infty$ 双模。
若位于
$(\mathfrak f^{(1)},\mathfrak f^{(0)})$
之上的滤过 $A_\infty$ 双模同态
$\varphi:D\to D'$ 满足两个
$\mathfrak f^{(i)}:C_i\to C'_i$
均严格，则它诱导上链映射
\[
\varphi_{0,0}:(D,\mathfrak n_{0,0})
\longrightarrow(D',\mathfrak n'_{0,0}).
\]
这里第 \textup{(2)} 项中
两个 $\mathfrak f^{(i)}$ 的严格性是假设的必要部分。
\end{statement}

\begin{proof}[子引理 5.2.20 的证明]
仍以 $\widehat d,\widehat d'$ 表示
$\mathfrak n,\mathfrak n'$ 诱导的余导子，
并以 $\widehat d^{(i)},\widehat d'^{(i)}$
表示 $\mathfrak m^{(i)},\mathfrak m'^{(i)}$
诱导的余导子。
若 $\mathfrak m_0^{(i)}=\mathfrak m_0'^{(i)}=0$，
则
$\widehat d^{(i)}(1)
=\widehat d'^{(i)}(1)=0$。

\textup{(1)} 对 $y\in D$，
\[
\begin{aligned}
\widehat d(y)
=\mathfrak n_{0,0}(y)
&+\widehat d^{(1)}(1)\otimes y\\
&+(-1)^{\deg' y}
y\otimes\widehat d^{(0)}(1).
\end{aligned}
\]
在上述假设下，
$\widehat d(y)=\mathfrak n_{0,0}(y)$；
由双模定义及 $\widehat d^2=0$，
即得第 \textup{(1)} 项。

\textup{(2)} 由于
$\widehat{\mathfrak f}^{(i)}(1)
=e^{\mathfrak f_0^{(i)}(1)}$，
而 $\mathfrak f^{(i)}$ 严格，
所以 $\widehat{\mathfrak f}^{(i)}(1)=1$。
由定义 3.7.7 与式 (3.7.2)，
\[
\begin{aligned}
\widehat d'\circ\widehat\varphi(y)
&=\widehat d'\bigl(
\widehat{\mathfrak f}^{(1)}(1)
\otimes\varphi_{0,0}(y)
\otimes\widehat{\mathfrak f}^{(0)}(1)
\bigr)\\
&=\widehat d'(\varphi_{0,0}(y))\\
&=\mathfrak n'_{0,0}\circ\varphi_{0,0}(y)
+\widehat d'^{(1)}(1)\otimes\varphi_{0,0}(y)\\
&\qquad+(-1)^{\deg' y}
\varphi_{0,0}(y)\otimes\widehat d'^{(0)}(1)\\
&=\mathfrak n'_{0,0}\circ\varphi_{0,0}(y).
\end{aligned}
\]
另一方面，
\[
\begin{aligned}
\widehat\varphi\circ\widehat d(y)
&=\widehat\varphi\bigl(
\mathfrak n_{0,0}(y)
+\widehat d^{(1)}(1)\otimes y
+(-1)^{\deg' y}y\otimes\widehat d^{(0)}(1)
\bigr)\\
&=\widehat\varphi(\mathfrak n_{0,0}(y))\\
&=\widehat{\mathfrak f}^{(1)}(1)
\otimes\varphi_{0,0}(\mathfrak n_{0,0}(y))
\otimes\widehat{\mathfrak f}^{(0)}(1)\\
&=\varphi_{0,0}\circ\mathfrak n_{0,0}(y).
\end{aligned}
\]
故由 $\widehat d'\circ\widehat\varphi
=\widehat\varphi\circ\widehat d$，
可得
$\mathfrak n'_{0,0}\circ\varphi_{0,0}
=\varphi_{0,0}\circ\mathfrak n_{0,0}$，
结论成立。
\end{proof}

子引理 5.2.20 在同伦含幺与弱无障碍情形下的版本
也可同法证明，
引理 5.2.19 因而得证。
\end{proof}
